\chapter{INTRODUCTION}

\section{Overview}
Autonomous mobile robotics and intelligent ground vehicles represent one of the most transformative frontiers in contemporary electrical, electronics, and computer engineering. From industrial Automated Guided Vehicles (AGVs) operating in structured warehouse facilities to full-scale autonomous road vehicles navigating complex urban intersections, the fundamental challenge remains consistent: developing robust, real-time sensing, perception, path planning, and motor actuation capabilities under stringent embedded computational and thermal constraints. 

Small-scale robotic platforms and autonomous micro-vehicles (often designed around radio-controlled chassis) serve as ideal cyber-physical testbeds for researching and validating self-driving architectures. Traditional embedded control systems implemented on micro-vehicles have largely depended on simple reactive mechanisms, such as infrared line-tracking sensors or localized sonar thresholds. While such reactive methods perform reliably in static, highly controlled environments, they inherently lack environmental semantic understanding, contextual awareness, and the adaptive intelligence required to negotiate dynamic obstacles, unmapped barricades, or multi-lane track intersections.

Recent breakthroughs in embedded Single-Board Computers (SBCs) and optimized deep learning runtimes have made it technically feasible to execute sophisticated convolutional neural networks and computer vision algorithms directly at the edge. By hosting edge intelligence on on-board embedded processors such as the Raspberry Pi, autonomous micro-vehicles can perform on-device inference without depending on fragile, high-latency remote cloud links. However, executing multi-stage perception pipelines—including high-resolution video capture, neural object detection, spatial corridor scanning, and closed-loop motor control—introduces severe resource contention. Ensuring low-latency execution, deterministic safety arbitration, and failsafe motor operation requires a tightly coupled, modular system architecture.

This project introduces a fully integrated, self-contained Edge-AI Autonomous RC Car architecture engineered for real-time vision-based navigation and obstacle avoidance. The platform integrates a Raspberry Pi single-board computer, an OV5647 Camera Serial Interface (CSI) sensor, dual ultrasonic distance sensors, a hardware Pulse Width Modulation (PWM) pan servo, an L298N dual H-bridge motor driver, and a comprehensive Flask telemetry web server. By combining optimized INT8-quantized object detection with spatial corridor clearance analysis and a deterministic safety watchdog, the vehicle demonstrates high reliability, responsive navigation, and robust collision avoidance in real-world operating conditions.

\section{Motivation and Objectives}
With the proliferation of autonomous robotics across smart logistics, disaster inspection, and smart transportation, developing efficient and self-sufficient embedded perception architectures is of critical importance. Many existing academic prototypes depend heavily on off-board desktop computers or high-end GPU workstations communicating via Wi-Fi to process camera video streams and transmit steering commands back to the vehicle. Such offloaded configurations suffer from network latency spikes, packet loss, bandwidth saturation, and absolute catastrophic failure whenever the wireless link drops. 

Conversely, attempting to run full-precision deep learning models locally on constrained single-board computers often results in severe frame rate drops (below 2--3 frames per second), excessive thermal throttling, and unacceptable reaction delays that compromise vehicle safety. Furthermore, generic object detection models trained on standard COCO datasets cannot inherently detect arbitrary floor obstacles, unclassified debris, or low-profile indoor barriers. 

The primary motivation of this project is to bridge these gaps by engineering a self-contained, deterministic, and fail-safe autonomous navigation architecture that operates completely offline on an embedded single-board computer while maintaining real-time responsiveness and rigorous safety contracts.

The key objectives of this project are:
\begin{enumerate}
    \item \textbf{Embedded Vision Pipeline:} To design and implement a low-latency, non-blocking camera capture and image preprocessing pipeline capable of extracting lane guidance centroids and topological junction features at 30 frames per second using native Picamera2 and OpenCV.
    \item \textbf{Edge-AI Semantic Perception:} To deploy a lightweight, 8-bit integer (INT8) quantized YOLOv8 object detection model on the Raspberry Pi using the LiteRT edge runtime, achieving deterministic inference latencies under 110~ms.
    \item \textbf{Spatial Corridor Clearance Analysis:} To develop a heuristic corridor scanner that partitions the immediate driving path into angular sectors, computing edge-density clearance scores to identify unobstructed evasion routes when primary lanes are blocked.
    \item \textbf{Multi-Sensor Fusion and Safety Arbitration:} To integrate ultrasonic range sensing with visual perception through an independent safety arbiter that overrides high-level planners, enforcing emergency braking within 15~cm and maintaining an active 0.6-second watchdog auto-stop timer.
    \item \textbf{Closed-Loop Motor and Actuation Control:} To implement proportional differential tank-mix motor control over software PWM at 100~Hz and zero-jitter camera pan survey using dedicated hardware PWM on GPIO pin 12.
    \item \textbf{Bidirectional Teleoperation and Telemetry:} To build an Apple-monochrome responsive web dashboard and Bluetooth RFCOMM interface delivering real-time MJPEG video feeds with Heads-Up Display (HUD) telemetry and manual override capabilities.
\end{enumerate}

\section{Report Organization}
This report is organized into five chapters. Chapter 1 introduces the project, outlining the domain background, motivation, core technical objectives, and system context. Chapter 2 provides an analytical review of existing literature on embedded computer vision, quantized edge deep learning, multi-sensor collision avoidance, and low-latency robotics telemetry frameworks. Chapter 3 presents the proposed system architecture, detailing the hardware pinout schematics, sensor fusion pipeline, software modular components, and the closed-loop execution workflow. Chapter 4 discusses the experimental results and performance benchmarks, including perception latency, corridor clearance accuracy, dynamic braking curves, telemetry stability, operational challenges, and future enhancements. Finally, Chapter 5 concludes the report with a summary of findings mapped to the design objectives and highlights key directions for future research. Standardized IEEE references, interface screenshots, and core executable Python source code excerpts are provided in the references and appendices.
